\section{Procedural generation}
\ac{PCG}, often reffered to as just \ac{PG}, is algorithmic generation of content. This content can take many forms including landscape, levels, textures or even 3D models. This leads to smaller file sizes since the program can generate the content using algorithms rather than having everything pre-made and saved. 

  This section will cover several methods of \ac{PG} and discuss similarities, differences and most importantly limitations of the methods. Additionally we will take a look at and what we are trying to achieve with the procedural generation in this project and propose a new method that will handle the problem of music generation in a new way. 

  \subsection{L-systems (Lindenmayer systems)}
  \textbf{\color{red}WIP section}\\
    Lindenmayer systems are a method of \ac{PG} based around the idea of creating rules for the generation which lead to self-similarity and fractal-like forms. This kind of \ac{PG} is often used in generating organic structures such as trees or other plant-like structures.

    If we observe Figure~\ref{fig:l_systems} we can see how two different L-systems can create both very rigid results in the form of the sierpinsky triangle and more organic structures like the fractal plant on the second line.

    \begin{figure}[h] % Sierpinksy and plant
      \begin{tabular}{ccc}

        \includegraphics[width=0.3\textwidth]{figures/04theory/sierpinski1.png} &
        \includegraphics[width=0.3\textwidth]{figures/04theory/sierpinski3.png} &
        \includegraphics[width=0.3\textwidth]{figures/04theory/sierpinski5.png} \\
        & (a) sierpinksi triangle & \\
      
        \includegraphics[width=0.3\textwidth]{figures/04theory/plant1.png} &
        \includegraphics[width=0.3\textwidth]{figures/04theory/plant3.png} &
        \includegraphics[width=0.3\textwidth]{figures/04theory/plant5.png} \\
        & (b) fractal plant & \\

        \includegraphics[width=0.3\textwidth]{figures/04theory/bush1.png} &
        \includegraphics[width=0.3\textwidth]{figures/04theory/bush3.png} &
        \includegraphics[width=0.3\textwidth]{figures/04theory/bush5.png} \\
        & (c) fractal bush & \\

      \end{tabular}

      \caption{Demonstration of the sierpinski triangle, a fractal plant and bush after one, three and five iterations repsectively.}\label{fig:l_systems}
    \end{figure}

    \begin{figure}[h] % Arrowhead sierpinski
      \begin{tabular}{ccc}

        \includegraphics[width=0.3\textwidth]{figures/04theory/arrowhead1.png} &
        \includegraphics[width=0.3\textwidth]{figures/04theory/arrowhead2.png} &
        \includegraphics[width=0.3\textwidth]{figures/04theory/arrowhead3.png} \\
        (a) first iteration & (b) second iteration & (c) third iteration \\
        \includegraphics[width=0.3\textwidth]{figures/04theory/arrowhead4.png} &
        \includegraphics[width=0.3\textwidth]{figures/04theory/arrowhead5.png} &
        \includegraphics[width=0.3\textwidth]{figures/04theory/arrowhead6.png} \\
        (d) fourth iteration & (e) fifth iterations & (f) sixth iteration \\

      \end{tabular}
      \caption{Approximation of the sierpinski triangle using the sierpinski arrowhead curve.}\label{fig:arrowhead_sierpinski}
    \end{figure}

  \subsection{PG using constraint satisfaction}
  \textbf{\color{red}WIP section}\\
    Similar to \ac{WFC} we have \ac{PG} using constraint satisfaction. This method uses given constraints for influencing the generation. It usually works by randomly generating something in the given domain, then iterates over the generated data and enforces constraints in order to make it adhere to a model designed by the programmer or end user. After several iteration over the domain we end up with a domain where the constraints are satisfied for all the objects in the domain.


    In Figure~\ref{fig:cons_adjacent} you can see how the constraints are set up for this example, and Figure~\ref{fig:domain_prop} shows the initial domain with how many different states are left for each cell in the domain as well as the domain after AC-3 propagation. Finally Figure~\ref{fig:cs_solution} shows the final domain that satisfies both the local and the global constraints. 

    \begin{figure}[h] % Nodes
      \begin{center}
        \includegraphics[width=0.95\textwidth]{figures/04theory/01_adjacency_constraints.pdf}
      \end{center}
      \caption{All different node types and how they are constrained. Each edge represents which note types can be adjacent in the domain.}\label{fig:cons_adjacent}
    \end{figure}

    \begin{figure}[h] % Initial domain and AC-3
      \begin{center}
        \includegraphics[width=0.95\textwidth]{figures/04theory/02_domain_propagation.pdf}
      \end{center}
      \caption{Initial domain and domain after AC-3 propagation. The AC-3 algoriths reduces the search space by ensuring arc consistency.}\label{fig:domain_prop}
    \end{figure}

    \begin{figure}[h] % Final map
      \begin{center}
        \includegraphics[width=0.6\textwidth]{figures/04theory/03_generated_solution.pdf}
      \end{center}
      \caption{Final generated map satisfying all constraints.}\label{fig:cs_solution}
    \end{figure}

  \subsection{Wave function collapse}
    \ac{WFC} is a method used to procedurally generate using either given constraints, or it can also use some input that it analyses and generates the constraints procedurally which are then used to generate something in the same style as the input. 

    It works by initializing the entire space in a superposition which means that the state is yet to be determined. We call this an \textit{unobserved} state, meaning we don't know what state it is in yet. These unobserved states are not complex numbers like they are in \ac{QM}, but the process is similar which is why it gets its name inspired by the field of \ac{QM}. Randomly assigning parts of the space we are able to collapse\footnote{collapsing a superposition means that we decide the actual state of the given position in the space} some of the \textit{wave function}\footnote{the wave function is simply the constraints what decide what the state of the object should be dependent on external factors}.
    
    The rules for how the tiling should propagate are derived from an example input. In Figure~\ref{fig:exemplar_and_rules} we see how from the simple input image the algorithm determines some local tiling rules for the different states and uses this as a basis for generating the new content. These generated rules are what make the output consistently similar and ensures we are outputting something that resembles the original input.

    \begin{figure}[h]
      \centering
        \includegraphics[width=0.95\textwidth]{figures/04theory/01_exemplar_and_rules.pdf}
      \caption{Given the exemplar on the left the \ac{WFC} algorithm learns the local tile rules shown on the right.}\label{fig:exemplar_and_rules}
    \end{figure}

    In Figure~\ref{fig:observe_propagate_steps} we can see how the number of possible states reduces as soon as a state is collapsed. This propagets outward and the process of assigning states continues until the entire domain has collapsed.

    \begin{figure}[h]
      \centering
        \begin{tabular}{cc}
          \includegraphics[width=0.4\textwidth]{figures/04theory/02_initial_wave.pdf} &
          \includegraphics[width=0.4\textwidth]{figures/04theory/02_after_first_collapse.pdf} \\
          (a) initial wave & (b) after first collapse \\
          \includegraphics[width=0.4\textwidth]{figures/04theory/02_after_propagation.pdf} &
          \includegraphics[width=0.4\textwidth]{figures/04theory/02_complete_collapse.pdf} \\
          (c) after propagation & (d) compete collapse
        \end{tabular}
      \caption{Shows how collapsing/observing a state propagates. An observation reduces the local uncertainty.}\label{fig:observe_propagate_steps}
    \end{figure}

    
    When the rules for the tiling have been established we can use them to generate new examples, in Figure~\ref{fig:generated_wfc_map} we are able to generate a larger domain using the same rules we got from the example image.

    \begin{figure}[h]
      \begin{center}
        \includegraphics[width=0.65\textwidth]{figures/04theory/03_generated_wfc_map.pdf}
      \end{center}
      \caption{An example output for a larger domain than the original exemplar. The consistent rules lead to consistently similar outputs.}\label{fig:generated_wfc_map}
    \end{figure}

    Finally we can see how the rules does not make it so that there is simply one output, we can change the seed for generation and end up with very different results, but all adhering to the initial rules derived from the exemplar in the very first step. 

    \begin{figure}[h]
      \begin{center}
        \includegraphics[width=0.95\textwidth]{figures/04theory/04_seed_variations.pdf}
      \end{center}
      \caption{Illustration that the learned rules stay consistend while the seed changes the output.}\label{fig:wfc_seed_variations}
    \end{figure}
  
  \subsection{Binary space partitioning}
  \textbf{\color{red}WIP section}\\
    \ac{BSP} can be used to generate map layouts like indoor environments, in the following figures you can see an example of how we can divide the entire domain into sections and used that to generate rooms.
  
    \begin{figure}[h] % Initial partitioning
      \begin{center}
        \includegraphics[width=0.5\textwidth]{figures/04theory/01_bsp_partitions.pdf}
      \end{center}
      \caption{Initial domain partitioning, made by recursively divinging the domain in two.}\label{fig:bsp_partitions}
    \end{figure}
    
    \begin{figure}[h] % Rooms and corridors
      \begin{center}
        \includegraphics[width=0.5\textwidth]{figures/04theory/02_rooms_and_corridors.pdf}
      \end{center}
      \caption{Generated rooms inside the partitions and connected them with corridors between rooms.}\label{fig:rooms_and_corridors}
    \end{figure}

    \begin{figure}[h]
      \begin{center}
        \includegraphics[width=0.5\textwidth]{figures/04theory/03_final_bsp_dungeon.pdf}
      \end{center}
      \caption{Rasterization of the generated dungeoun to show "final product".}\label{fig:final_bsp_dungeoun}
    \end{figure}

  \subsection{Random noise and seeds}
  \textbf{\color{red}WIP section}\\
    Perlin noise is a very common way to generate structures like worlds etc... It is very good at creating smooth transitions from one state to another, like height maps.
    This is demonstrated in the following figures below with some different techniques.
    
    \begin{figure}[h]
      \begin{center}
        \includegraphics[width=0.95\textwidth]{figures/04theory/01_seed_variation.pdf}
      \end{center}
      \caption{}\label{fig:seed_variations}
    \end{figure}
    \begin{figure}[h]
      \begin{center}
        \includegraphics[width=0.95\textwidth]{figures/04theory/02_noise_octaves.pdf}
      \end{center}
      \caption{}\label{fig:noise_octaves}
    \end{figure}
    \begin{figure}[h]
      \begin{center}
        \includegraphics[width=0.95\textwidth]{figures/04theory/03_noise_to_terrain.pdf}
      \end{center}
      \caption{}\label{fig:noise_to_terrain}
    \end{figure}
    \begin{figure}[h]
      \begin{center}
        \includegraphics[width=0.95\textwidth]{figures/04theory/04_parameter_control.pdf}
      \end{center}
      \caption{}\label{fig:parameter_control}
    \end{figure}

  \subsection{Random walks}
    Random walk procedural generation method that can be used to organic looking structures like maps and textures. It can also be used to describe natural phenomena like bacterial movement. It has been applied in financial economics and computer science

    The random walk works by having an entity walk around randomly. In some cases the random walk is also called the \textit{drunkard's walk}. Then the path is traced, either to create a circumference of an area or the traversed spaces become the area itself. Oftentimes the drunk walker has a set amount of time or steps it can walk before it is terminated, then a new walker is placed at a random location, this helps to smooth out the generated area.

    The same technique can be used to create textures as well, in Figure~\ref{fig:random_walks} we can see how changing the colours and directional bias we can generate textures that imitate real life. In the tree and sand example the directional bias is vertical for the tree texture and diagonal for the sand texture. For the map texture it is uniform directional bias, but colouring to indicate water, land and mountains. The least traveled to square get more blue, and the more frequent it goes to green and brown/white to indicate mountains.

    \begin{figure}[h]
      \centering
        \includegraphics[width=0.32\textwidth]{figures/04theory/ex_wood.png}
        \includegraphics[width=0.32\textwidth]{figures/04theory/ex_sand.png}
        \includegraphics[width=0.32\textwidth]{figures/04theory/ex_map.png}
      \caption{Three different textures generated using random walks.}\label{fig:random_walks}
    \end{figure}

  \subsection{Discussion}
  \textbf{\color{red}WIP section}\\
    Now that we have taken a look into a few methods of \ac{PG} we observe that there are many similarities between the methods. Common for all forms of \ac{PG} is that they utilize some form of randomness. The randomness in a core part of procedural generation, if there was no randomness we would generate the same output every time and that would be the same as just creating the content yourself.
